\cotefile{7}
\scanpage{61}
\[
H^*((U,X'),F)\simeq H^*(X',\unsure{F_{(U,X')}})
\]
Supposons pour \unsure{simplifier} $X$ affine, \deleted{\unsure{et affine sur $S$}}, donc les $\Delta^{(\nu)}(i)$ affines. \unsure{Supposons} que les $F_{\Delta^{(\nu)}(i)}$ soient \uline{quasi-cohérents}\inserted{\uline{et s'induisent mutuellement}}. Alors \unsure{on trouve donc}
\[
H^q(\widetilde{\Delta_X^{(\nu)}(i)},F)=0\qquad\text{si }q>0
\]
\inserted{$H^0(\widetilde{\Delta_X^{(\nu)}(i)},F)$ pour $i$ variable est un \unsure{syst. proj. strict},}
d'où par les \unsure{arguments à la} M.L.
\[
H^q(\widetilde X^{(\nu+1)},F)
=H^q(\varinjlim_i\widetilde{\Delta_X^{(\nu)}(i)},F)
=\varprojlim_i H^q(\widetilde{\Delta_X^{(\nu)}(i)},F)
\]
donc
\[
\left\{\begin{aligned}
H^q(\widetilde X^{(\nu+1)},F)&=0&&\text{si }q>0\\
H^0(\widetilde X^{(\nu+1)},F)&=\varprojlim_i F(\Delta_X^{(\nu)}(i))
\end{aligned}\right.
\]
et par suite un isom \unsure{can}
\[
\left\{\begin{gathered}
H^*(X_{\mathrm{strat}},F)\simeq
H^*\bigl(\nu\longmapsto F(\underbrace{\unsure{\widehat{X^{(\nu+1)}}_{/\deleted{\Delta}X}}}_{\text{complété formel}})\bigr)\\
\text{si $X$ affine, les $F_{\Delta_X^{(\nu)}(i)}$ quasi-cohérents}\quad
\left(\text{p. ex. }F=\unsure{\mathcal O_{X_{\mathrm{strat}}}}\right)
\end{gathered}\right.
\]
Dans le cas où on suppose seulement que les $F_{\Delta_X^{(\nu)}(i)}$ sont qu. coh.\deleted{\unsure{(quasi-induits\ldots ?)}}\inserted{\uline{et s'induisent mutuellement}}, mais $X$ pas \unsure{néc.} affine, on \unsure{voit} que pour tt entier $\nu$, on a un
\mspage{60}
